\documentclass[11pt,a4paper]{article}

% ----- Premium Typography & Formatting -----
\usepackage[utf8]{inputenc}
\usepackage[T1]{fontenc}
\usepackage{mathpazo} % Elegant Palatino font for text and math
\usepackage{microtype} % Improves text justification
\usepackage{amsmath, amssymb}
\usepackage{graphicx}
\usepackage{booktabs}
\usepackage{geometry}
\usepackage{xcolor}
\usepackage{caption}
\usepackage[most]{tcolorbox} % For beautiful highlight boxes
\usepackage{hyperref}
\usepackage{enumitem}

\geometry{top=1in, bottom=1.2in, left=1in, right=1in}

% ----- Custom Colors -----
\definecolor{primary}{RGB}{44, 62, 80}
\definecolor{secondary}{RGB}{52, 152, 219}
\definecolor{highlight}{RGB}{231, 76, 60}

\hypersetup{
    colorlinks=true,
    linkcolor=secondary,
    filecolor=secondary,      
    urlcolor=secondary,
    citecolor=primary,
}

% ----- Standard Title Section (Error-Free) -----
\title{\vspace{-2cm}\textbf{\color{primary} Solving the Exploding Gradient Problem in BPTT} \\ \Large \color{secondary} Empirical Analysis of Gradient Norm Clipping in Deep Recurrent Networks}
% Blank date

\begin{document}

\maketitle

% ----- Abstract -----
\begin{tcolorbox}[colback=gray!5, colframe=primary!80, title=\textbf{Abstract}, boxrule=0.8pt, arc=3pt]
\textit{This report empirically investigates the exploding gradient phenomenon inherent in Deep Recurrent Neural Networks (RNNs) trained via Backpropagation Through Time (BPTT). Over a series of rigorous ablation studies, we isolated the structural vulnerabilities of deep temporal unrolling. We demonstrate the mathematical superiority of \textbf{Gradient Norm Clipping} and \textbf{Constant Error Carousels (LSTMs)} in stabilizing high-variance optimization landscapes.}
\end{tcolorbox}

\vspace{0.5cm}

% ----- Section 1 -----
\section{Introduction and Theoretical Background}
The exploding gradient problem is a fundamental challenge when training deep neural networks on long sequences. When applying the chain rule across sequence length $T$ to compute the gradient of the loss $\mathcal{L}$ with respect to the initial hidden state $h_0$:

\begin{tcolorbox}[colback=blue!5, colframe=blue!30, boxrule=0.5pt, arc=2pt]
\begin{equation}
\frac{\partial \mathcal{L}}{\partial h_0} = \frac{\partial \mathcal{L}}{\partial h_T} \prod_{t=1}^{T} \frac{\partial h_t}{\partial h_{t-1}}
\end{equation}
\end{tcolorbox}

The magnitude of this gradient is heavily dominated by the eigenvalues of the hidden-to-hidden weight matrix. If the largest singular value $\lambda_{max} > 1$, the gradient norm grows exponentially with each multiplication, causing violent updates in parameter space.

\section{Methodology}
Our methodology relies on a controlled, synthetic time-series dataset featuring dual-frequency superimposed sine waves embedded with Gaussian noise. 
\begin{itemize}[label=\textcolor{secondary}{\textbullet}]
    \item \textbf{Sequence Length ($T=60$):} Forces the network to backpropagate across 60 dense matrix multiplications.
    \item \textbf{Initialization:} We utilized strict \textit{Orthogonal Weight Initialization}, ensuring the initial eigenvalues start at exactly $\lambda = 1.0$.
\end{itemize}

\section{Experimental Findings}

\subsection{Experiment 1: The Anatomy of an Explosion}
Without intervention, the BPTT Jacobian chain rule scales eigenvalues aggressively. Our unclipped baseline experienced severe instability. Pre-clip gradient norms spiked up to $\mathbf{35.82}$. These massive $L_2$ norm deviations correspond directly to violent spikes in validation loss, immediately undoing the optimizer's convergence trajectory.

\begin{figure}[ht]
    \centering
    % Images are active now!
    \includegraphics[width=1\textwidth]{exp1.png}
    \caption{\textit{Training instability without gradient clipping. Notice the massive gradient spikes corresponding with loss degradation.}}
    \label{fig:exp1}
\end{figure}

\subsection{Experiment 2: Establishing Control via Norm Clipping}
By implementing gradient norm clipping, we scale the gradient vector $g$ conditionally:
\begin{equation}
g_{new} = g \times \frac{c}{\Vert g \Vert_2}
\end{equation}
Because it acts as a universal scalar, it enforces a strict boundary ($C=1.0$) on the optimizer step size without distorting the multi-dimensional direction of steepest descent.

\subsection{Experiment 3: Threshold Sensitivity}
An ablation study across $C \in \{0.1, 0.5, 1.0, 5.0\}$ revealed that the threshold $C$ is highly sensitive:
\begin{itemize}[label=\textcolor{highlight}{\textbullet}]
    \item \textbf{$C=0.1$:} Triggered on every single iteration (360/360 times). The optimizer was choked.
    \item \textbf{$C=5.0$:} Triggered only 33 times out of 360. Maximum raw norms still breached $\mathbf{75.80}$, offering minimal protection.
    \item \textbf{$C=1.0$:} The optimal threshold, cleanly truncating exponential spikes while allowing native gradient flows.
\end{itemize}

\subsection{Experiment 4: The Depth Penalty}
We analyzed variance across $L \in \{2, 4, 6\}$ hidden layers. The gradient norm standard deviation ($\sigma$) increased from \textbf{1.72} (2-layer) to \textbf{5.46} (6-layer) --- a staggering $>3\times$ increase in variance. Deeper networks exponentially compound the Jacobian matrix multiplications across both spatial and temporal axes.

\begin{figure}[ht]
    \centering
    % Images are active now!
    \includegraphics[width=0.8\textwidth]{exp4_violin.png}
    \caption{\textit{Violin plot demonstrating the geometric increase in gradient variance as network depth increases.}}
    \label{fig:exp4}
\end{figure}

\subsection{Bonus Experiment: Architectural Superiority}
Advanced gating mechanisms (LSTMs) utilize additive cell states. When the forget gate is saturated, the derivative approaches $1.0$. Both LSTM and GRU variations massively outperformed Vanilla RNNs, establishing smooth, deeply converged validation losses immune to exponential decay.

\vspace{0.5cm}

% ----- Production Recommendations -----
\begin{tcolorbox}[colback=secondary!5, colframe=secondary!80, title=\textbf{Production Recommendations}, boxrule=0.8pt, arc=3pt]
For deployment in production sequence-modeling environments, we mandate the following best practices:
\begin{enumerate}
    \item \textbf{Base Architecture:} LSTM or GRU to structurally eliminate vanishing gradients.
    \item \textbf{Algorithmic Safety Net:} Gradient norm clipping with $C \in [0.5, 1.0]$.
    \item \textbf{Weight Initialization:} Orthogonal initialization for all recurrent matrices.
    \item \textbf{Optimization:} AdamW optimizer with Cosine Annealing Learning Rate scheduling.
    \item \textbf{Robust Engineering:} Strict NaN detection during the backward pass for early-stopping.
\end{enumerate}
\end{tcolorbox}

\end{document}
